\documentclass[11pt]{article}
\usepackage[a4paper,margin=25mm]{geometry}
\usepackage{amsmath,amssymb,amsthm}
\usepackage[hidelinks]{hyperref}
\newtheorem{theorem}{Theorem}
\newtheorem{lemma}{Lemma}
\title{The Wrong Spectral Scale for Quadratic-Residue Graphs:\\A Disproof of Conjecture 00000001071}
\author{GitHub: gaochengzhecpu}
\date{}
\begin{document}
\maketitle
\begin{abstract}
The source compares the smallest adjacency eigenvalue of a quadratic-residue
graph with $-2\sqrt p$ and claims an error at most $p^{-1/4+\varepsilon}$.
For the standard Paley graph, the discrepancy is at least
$(3\sqrt p-1)/2$. If the source's parenthetical ``adjacency by $\pm1$''
instead means differences $\pm1$, its graph is a cycle and the discrepancy
is at least $2\sqrt p-2$. Both lower bounds are unbounded along the
appropriate prime orders. The graph definitions, character identities,
matrix arguments, prime-order quantifiers and negation of the error estimate
are all formalized in Lean.
\end{abstract}

\paragraph{Reading of the source.}
Both language versions name a quadratic-residue graph and also give the
parenthetical adjacency rule $\pm1$. These descriptions need not describe
the same graph. We therefore treat both explicitly. The standard undirected
Paley graph $P_p$, for a prime $p\equiv1\pmod4$, has vertex set
$\mathbb F_p$, with distinct vertices adjacent when their difference is a
square. Under the literal difference-$\pm1$ rule, the graph is the cycle
$C_p$ on $\mathbb Z/p\mathbb Z$. We use the ordinary, unnormalized adjacency
matrix in each case. The phrase ``deviation from $-2\sqrt p$'' is read as
$|\lambda_{\min}+2\sqrt p|$. Allowing a constant and an arbitrary starting
point makes the estimate weaker than its displayed constant-one form, and
our disproof covers this weaker asymptotic interpretation as well.

\begin{theorem}\label{thm:main}
Let $\lambda_{\min}(G)$ be the smallest eigenvalue of the real adjacency
matrix of a finite graph $G$.
\begin{enumerate}
\item For every prime $p\equiv1\pmod4$,
\[
 |\lambda_{\min}(P_p)+2\sqrt p|\ \geq\ \frac{3\sqrt p-1}{2}.
\]
\item For every prime $p\geq3$,
\[
 |\lambda_{\min}(C_p)+2\sqrt p|\ \geq\ 2\sqrt p-2.
\]
\end{enumerate}
For either family, for every real $B$ and every integer cutoff $N$, an
admissible prime $p\geq N$ has discrepancy greater than $B$.
Consequently neither family satisfies
$|\lambda_{\min}+2\sqrt p|=O(p^{-1/4+\varepsilon})$ for every
$\varepsilon>0$; already $\varepsilon=1/8$ fails.
\end{theorem}

\paragraph{Quadratic-character identities.}
We give the argument over any finite field $F$ of odd characteristic,
writing $q=|F|$ and assuming that $-1$ is a square. Let
$\chi:F\longrightarrow\{0,1,-1\}$ be its quadratic character, extended by
$\chi(0)=0$. Thus $\chi(-1)=1$, $\chi(u)^2=1$ for $u\ne0$, and
$\sum_{u\in F}\chi(u)=0$.

First,
\begin{equation}\label{eq:jacobi}
 \sum_{t\in F}\chi(t)\chi(1-t)=-1.
\end{equation}
Indeed the quadratic character equals its inverse as a multiplicative
character, so the summand equals $\chi(t/(1-t))$ for $t\ne1$; at $t=1$
the original summand is zero. The map
\[
 F\setminus\{1\}\longrightarrow F\setminus\{-1\},\qquad
 t\longmapsto\frac{t}{1-t}
\]
is a bijection with inverse $u\mapsto u/(1+u)$. The sum is therefore
$\sum_{u\ne-1}\chi(u)=-\chi(-1)=-1$.

For $a\ne b$, put $d=a-b$ and substitute $x=b+dt$ to obtain
\begin{equation}\label{eq:correlation}
 \sum_{x\in F}\chi(a-x)\chi(x-b)
 =\chi(d)^2\sum_{t\in F}\chi(1-t)\chi(t)=-1.
\end{equation}
For $a=b$, the summand is $\chi(a-x)^2$, since $\chi(-1)=1$, and its sum
is $q-1$.

\paragraph{The actual Paley adjacency matrix.}
Index the following real matrices by $F$. Let $J$ be the all-ones matrix,
let $I$ be the identity, and set
\[
 Q_{ab}=\chi(a-b),\qquad A=\frac{J-I+Q}{2}.
\]
On the diagonal $A_{aa}=0$; off the diagonal it is one exactly when
$a-b$ is a square, and otherwise zero. Thus $A$ is precisely the adjacency
matrix of the standard Paley graph, rather than a surrogate matrix.
It is symmetric because $\chi(-1)=1$.
The diagonal and off-diagonal calculations above give
\begin{equation}\label{eq:matrix}
 Q^2=qI-J.
\end{equation}
Each column of $Q$ sums to zero, by translating the character sum.
Consequently every column of $A$ sums to $k=(q-1)/2$.

Let $Ax=\lambda x$ with $x\ne0$, and put $s=\sum_a x_a$.
Summing the eigenvector equation gives $ks=\lambda s$. If $\lambda=k$,
then $\lambda\geq0$. Otherwise $s=0$, hence $Jx=0$ and
\[
 Qx=(2\lambda+1)x,\qquad
 Q^2x=(2\lambda+1)^2x=qx.
\]
A nonzero coordinate of $x$ gives $(2\lambda+1)^2=q$. In either case,
\begin{equation}\label{eq:lower}
 \lambda\geq-\frac{1+\sqrt q}{2}.
\end{equation}
Applying this to the smallest adjacency eigenvalue and using $|u|\geq u$
yields
\[
 |\lambda_{\min}+2\sqrt q|\geq\lambda_{\min}+2\sqrt q
 \geq\frac{3\sqrt q-1}{2}.
\]
This proof needs only a spectral lower bound, not an unproved claim about
the occurrence or multiplicities of either nontrivial eigenvalue.

\paragraph{The literal difference-$\pm1$ graph.}
More generally, let $G$ have maximum degree at most two and $Ax=\lambda x$
with $x\ne0$. Choose $v$ maximizing $|x_v|$, which is positive. Then
\[
 |\lambda|\,|x_v|
 =\left|\sum_{w\sim v}x_w\right|
 \leq\sum_{w\sim v}|x_w|\leq2|x_v|.
\]
Thus $|\lambda|\leq2$. For $C_p$ this gives
$\lambda_{\min}\geq-2$ and the second lower bound in
Theorem~\ref{thm:main}.

\paragraph{Prime orders and the asymptotic quantifiers.}
There are arbitrarily large primes, and also arbitrarily large primes
$p\equiv1\pmod4$. For the latter elementary fact, suppose that
$r_1,\ldots,r_m$ list all such primes and take
$M=(2r_1\cdots r_m)^2+1$. Every prime divisor $\ell$ of $M$ is odd,
and $-1$ is a square modulo $\ell$. The corresponding nonzero residue
has order four, so $4\mid\ell-1$. Also $\ell$ is not one of the $r_i$,
a contradiction. These infinitude statements imply unboundedness.

Given $B\in\mathbb R$ and $N$, take an admissible prime greater than both
$N$ and $(|B|+2)^2$. Then $\sqrt p>|B|+2$, and each applicable lower bound
in Theorem~\ref{thm:main} exceeds $B$.
If the source's error estimate held at $\varepsilon=1/8$, some $C\geq0$
and $N$ would give discrepancy at most $Cp^{-1/8}\leq C$ for all
admissible primes $p\geq N$. This contradicts the preceding result with
$B=C$. Failure of the error clause already disproves the conjunction in
the source; no separate claim about its spectral-gap clause is needed.

\paragraph{Formal verification and scope.}
The file \texttt{lean/Main.lean} formalizes both readings. The standard
Paley graph is a Mathlib \texttt{SimpleGraph} over \texttt{ZMod p}.
Its adjacency is proved equivalent to distinctness and the standard
\texttt{IsSquare} predicate on the difference. Its matrix is Mathlib's
adjacency matrix. The quadratic character is Mathlib's actual finite-field
character. The Jacobi sum identity is the library theorem for a character
and its inverse; the affine change of variables, character correlation,
matrix identity, column sums and spectral lower bound are proved in the
submission. No desired spectrum is assumed.

The smallest eigenvalue is the minimum of Mathlib's complete real
eigenvalue list for a Hermitian matrix. A member of that list is supplied
with a nonzero vector from its orthonormal eigenvector basis. The cycle
uses Mathlib's \texttt{cycleGraph}, including its degree theorem. The
formal prime statements have arbitrary real bounds and arbitrary natural
cutoffs. They prove the negations of universally quantified error claims,
including arbitrary nonnegative error constants. The final theorems are
\begin{center}\small
\texttt{Conjecture1071.Paley.not\_spectralClaim}\\
\texttt{Conjecture1071.not\_spectralClaim}.
\end{center}
All mathematical bridges used above are covered by Lean, except that the
elementary exposition of infinitude is replaced there by Mathlib's proved
infinitude theorems. The argument does not concern an unspecified different
graph family or an unstated rescaling of its adjacency matrix.

\paragraph{Source and reproduction.}
The original bilingual source is preserved in \texttt{SOURCE.md}; its
commit and SHA-256 digest are recorded in
\texttt{verification/SOURCE\_PROVENANCE.md}. The Lean project pins Lean
4.19.0 and Mathlib commit
\begin{center}\small\texttt{c44e0c8ee63ca166450922a373c7409c5d26b00b}.\end{center}
From the \texttt{lean} directory run
\begin{center}\small\texttt{lake env lean -DwarningAsError=true Main.lean}.\end{center}
The printed axiom reports use only \texttt{propext},
\texttt{Classical.choice} and \texttt{Quot.sound}.
This submission was developed by a delegated AI agent with a self-review;
the coordinating agent reviews it separately. There is no claim of an
external independent review.
\end{document}
